\chapter{Tier 0: the learned superelement}

At tier 0 a learned expert is not asked for a port matrix.  It returns
\emph{fields}: for each pattern of boundary values, an interior field to go
with it.  The host computes the port matrix itself, as the energy Gram matrix
of those fields, assembles one interface system, and solves it once.  There is
no iteration.

The statements of this chapter hold for \emph{every} output of the network,
however wrong.  They are the tier-0 guarantee of the architecture document
(its Theorems 5 and 6, Proposition 7 and Corollary 11): the learned system is
never indefinite, it is solvable whenever the classical one is, its answer is
the best its modes can represent, and its sensitivity is governed by the
classical problem's own constant.

Everything is finite-dimensional linear algebra.  The statements are written
for real vector spaces and bilinear forms, so that nothing depends on a choice
of basis for the fields; with coordinates they are the matrix statements of the
document.

\section{T25 (i) to (iii): the Gram port matrix}

\paragraph{In plain words.}
A piece of the domain has unknowns inside it and unknowns on its boundary
faces, and an energy that is never negative.  A port mode is a pattern of
boundary values; its exact constraint mode is the interior field the physics
pairs with it, and the exact port matrix is the energy Gram matrix of those
pairs.  Put any other interior fields in place of the exact ones.  Then the
Gram matrix is still symmetric and positive semidefinite; it exceeds the exact
one by exactly the Gram matrix of the field errors in the interior energy; and
after assembly over all pieces the learned interface matrix dominates the exact
one, so it inherits every lower bound the exact one has.  An error in the
fields can make a piece stiffer, never softer, and never indefinite.

\begin{definition}[The energy of a piece]\label{def:piece}
  \lean{Atlas.Piece, Atlas.Piece.M, Atlas.Piece.IsExactMode}
  \leanok
  Let $U$ (cell values) and $F$ (face values) be real vector spaces.  A
  \emph{piece} is three bilinear forms, written with matrices as
  $a(u,v)=u^\top A_Iv$ on $U$, $b(u,\mu)=u^\top B\mu$ on $U\times F$ and
  $d(\lambda,\mu)=\lambda^\top D\mu$ on $F$, with $a$ and $d$ symmetric and
  \[
    a(u,u)-2\,b(u,\lambda)+d(\lambda,\lambda)\ \ge\ 0
    \qquad\text{for all }u\in U,\ \lambda\in F .
  \]
  Its energy form on pairs is
  \[
    M\bigl((u,\lambda),(v,\mu)\bigr)=a(u,v)-b(u,\mu)-b(v,\lambda)+d(\lambda,\mu),
    \qquad M=\begin{pmatrix}A_I&-B\\-B^\top&D\end{pmatrix}.
  \]
  An interior field $h\in U$ is the \emph{exact constraint mode} of the
  boundary pattern $q\in F$ if $a(v,h)=b(v,q)$ for every $v\in U$, that is,
  $A_Ih=Bq$.
\end{definition}

\begin{definition}[The Gram port matrix]\label{def:gram}
  \lean{Atlas.Piece.gram}
  \leanok
  \uses{def:piece}
  For interior fields $h_k\in U$ and boundary patterns $q_k\in F$, indexed by a
  finite set, the Gram port matrix has entries
  \[
    \Lambda[h]_{kl}=M\bigl((h_k,q_k),(h_l,q_l)\bigr).
  \]
  With the exact constraint modes $H_k$ it is the exact port matrix
  $\Lambda=\Lambda[H]$; with a network's fields $\hat H_k$ it is the matrix
  $\tilde\Lambda=\Lambda[\hat H]$ the host computes.
\end{definition}

\begin{theorem}[T25 (i), positive semidefinite for every output]\label{thm:T25-i}
  \lean{Atlas.Piece.gram_posSemidef}
  \leanok
  \uses{def:gram}
  For any fields $h_k$ and patterns $q_k$, $\Lambda[h]$ is symmetric and
  $x^\top\Lambda[h]\,x\ge0$ for every $x$.
\end{theorem}

\begin{proof}
  \leanok
  $\Lambda[h]$ is symmetric because $M$ is.  With
  $w=\bigl(\sum_kx_kh_k,\ \sum_kx_kq_k\bigr)$, bilinearity gives
  $x^\top\Lambda[h]\,x=M(w,w)\ge0$.
\end{proof}

\begin{proposition}[T25 (i), the definite case]\label{prop:T25-i-def}
  \lean{Atlas.Piece.gram_posDef}
  \leanok
  \uses{def:gram}
  If $M(w,w)=0$ only for $w=0$, and the patterns $q_k$ are linearly
  independent, then $\Lambda[h]$ is positive definite for any fields $h_k$.
\end{proposition}

\begin{proof}
  \leanok
  \uses{thm:T25-i}
  If $x^\top\Lambda[h]\,x=0$ then $w=0$ in the proof of
  Theorem~\ref{thm:T25-i}; its face part is $\sum_kx_kq_k=0$, so $x=0$.
\end{proof}

\begin{theorem}[T25 (ii), the error identity]\label{thm:T25-ii}
  \lean{Atlas.Piece.gram_sub_gram, Atlas.Piece.quadratic_gram_sub}
  \leanok
  \uses{def:gram}
  Let $H_k$ be the exact constraint modes of the patterns $q_k$, let $h_k$ be
  any fields, and $E_k=h_k-H_k$.  Then
  \[
    \Lambda[h]_{kl}-\Lambda[H]_{kl}=a(E_k,E_l)\qquad\text{for all }k,l,
  \]
  and so, for every coordinate vector $x$,
  \[
    x^\top\bigl(\Lambda[h]-\Lambda[H]\bigr)x
      =a\Bigl(\sum_kx_kE_k,\ \sum_kx_kE_k\Bigr).
  \]
  In matrices: $\tilde\Lambda-\Lambda=E^\top A_IE$.
\end{theorem}

\begin{proof}
  \leanok
  Write $(h_k,q_k)=(H_k,q_k)+(E_k,0)$ and expand $M$ in both arguments.  The
  cross terms are $M\bigl((E_k,0),(H_l,q_l)\bigr)=a(E_k,H_l)-b(E_k,q_l)=0$ by
  the definition of $H_l$, and the same with $k$ and $l$ exchanged.  The last
  term is $M\bigl((E_k,0),(E_l,0)\bigr)=a(E_k,E_l)$.
\end{proof}

\begin{corollary}[T25 (ii), consequences]\label{cor:T25-ii}
  \lean{Atlas.Piece.gram_sub_gram_posSemidef, Atlas.Piece.trace_gram_sub,
    Atlas.Piece.quadratic_gram_sub_le}
  \leanok
  \uses{thm:T25-ii}
  With the notation of Theorem~\ref{thm:T25-ii}:
  \begin{enumerate}
    \item $\Lambda[h]-\Lambda[H]$ is positive semidefinite;
    \item $\operatorname{tr}\Lambda[h]-\operatorname{tr}\Lambda[H]=\sum_ka(E_k,E_k)$;
    \item if $s:U\to\R$ is any measure of size with $a(u,u)\le\alpha\,s(u)^2$
      for all $u$, $\alpha\ge0$, and
      $s\bigl(\sum_kx_kE_k\bigr)^2\le\varepsilon^2\,|x|^2$ for all $x$, then
      $x^\top\bigl(\Lambda[h]-\Lambda[H]\bigr)x\le\alpha\,\varepsilon^2\,|x|^2$.
      With Euclidean norms this is
      $\norm{\tilde\Lambda-\Lambda}_2\le\norm{A_I}_2\norm{E}_2^2$: the error of
      the port matrix is quadratic in the error of the fields.
  \end{enumerate}
\end{corollary}

\begin{proof}
  \leanok
  \uses{thm:T25-ii}
  The interior energy $a(u,u)=M\bigl((u,0),(u,0)\bigr)$ is non-negative, which
  gives 1 from the quadratic form of Theorem~\ref{thm:T25-ii}.  Item 2 is the
  diagonal of the identity.  Item 3 is the quadratic form with the two bounds.
\end{proof}

\begin{definition}[Assembly]\label{def:assemble}
  \lean{Atlas.assemble}
  \leanok
  For pieces $i$ in a finite set, port matrices $L_i$ and matrices $R_i$ that
  read piece $i$'s port coordinates off the global coordinate vector, the
  assembled interface matrix is $\sum_iR_i^\top L_iR_i$.
\end{definition}

\begin{theorem}[T25 (iii), the assembled matrix]\label{thm:T25-iii}
  \lean{Atlas.assemble_gram_sub_posSemidef, Atlas.assemble_gram_coercive,
    Atlas.assemble_gram_posDef, Atlas.assemble_gram_existsUnique}
  \leanok
  \uses{def:assemble, cor:T25-ii}
  Let $\mathsf S=\sum_iR_i^\top\Lambda_i[H_i]R_i$ be assembled from the exact
  constraint modes of every piece and
  $\tilde{\mathsf S}=\sum_iR_i^\top\Lambda_i[h_i]R_i$ from any fields.  Then
  \begin{enumerate}
    \item $\tilde{\mathsf S}-\mathsf S$ is positive semidefinite;
    \item if $\beta\,|c|^2\le c^\top\mathsf Sc$ for all $c$, then
      $\beta\,|c|^2\le c^\top\tilde{\mathsf S}c$ for all $c$ (with the largest
      such $\beta$: $\lambda_{\min}(\tilde{\mathsf S})\ge\lambda_{\min}(\mathsf S)$);
    \item if $\mathsf S$ is positive definite, so is $\tilde{\mathsf S}$, and
      $\tilde{\mathsf S}c=\chi$ has exactly one solution for every $\chi$.
  \end{enumerate}
\end{theorem}

\begin{proof}
  \leanok
  \uses{cor:T25-ii}
  $\tilde{\mathsf S}-\mathsf S=\sum_iR_i^\top(\Lambda_i[h_i]-\Lambda_i[H_i])R_i$
  is a sum of matrices $R^\top PR$ with $P$ positive semidefinite by
  Corollary~\ref{cor:T25-ii}.  Then
  $c^\top\tilde{\mathsf S}c=c^\top\mathsf Sc+c^\top(\tilde{\mathsf S}-\mathsf S)c
  \ge c^\top\mathsf Sc$, which gives 2 and the definiteness in 3.  A positive
  definite matrix is invertible.
\end{proof}

\section{T25: energy optimality}

\paragraph{In plain words.}
The undivided discrete solution is the admissible field of least total energy.
The host's interface system, formed from the network's fields alone, says
exactly that the rebuilt field has the least energy among the fields that can
be rebuilt from the network's modes.  For a symmetric problem, least energy
over a smaller set is the same as closest to the true minimiser in the energy
size.  So the coupled answer is the best approximation of the undivided
solution that the network's modes can represent: the coupling adds no error of
its own.  Its error is at most the truncation error of the classical
superelement plus the energy of the network's field errors.

\begin{definition}[A quadratic energy]\label{def:quad}
  \lean{Atlas.QuadEnergy, Atlas.QuadEnergy.energy, Atlas.QuadEnergy.norm}
  \leanok
  On a real vector space $X$, let $m$ be a symmetric bilinear form with
  $m(w,w)\ge0$ and $\ell$ a linear functional.  The energy is
  $\mathcal E(w)=\tfrac12m(w,w)-\ell(w)$ and the energy size is
  $\norm{w}_m=\sqrt{m(w,w)}$, a seminorm.
\end{definition}

\begin{lemma}[Minimisers are the stationary points]\label{lem:stationary}
  \lean{Atlas.QuadEnergy.stationary_of_isMin, Atlas.QuadEnergy.isMin_of_stationary,
    Atlas.QuadEnergy.energy_eq_add}
  \leanok
  \uses{def:quad}
  Let $U_0\subseteq X$ be a subspace and $u\in X$.  Then
  $\mathcal E(u)\le\mathcal E(u+v)$ for all $v\in U_0$ if and only if
  $m(u,v)=\ell(v)$ for all $v\in U_0$.  In that case, for every $w$ with
  $w-u\in U_0$,
  \[
    \mathcal E(w)=\mathcal E(u)+\tfrac12\norm{w-u}_m^2 .
  \]
\end{lemma}

\begin{proof}
  \leanok
  $\mathcal E(u+tv)-\mathcal E(u)=t\,\bigl(m(u,v)-\ell(v)\bigr)+\tfrac12t^2m(v,v)$.
  If this is non-negative for all real $t$, the coefficient of $t$ is zero.
  Conversely, if that coefficient is zero the difference is
  $\tfrac12t^2m(v,v)\ge0$; $t=1$ and $v=w-u$ give the identity.
\end{proof}

\begin{lemma}[The abstract lemma: C\'ea with constant one]\label{lem:cea}
  \lean{Atlas.QuadEnergy.closest_iff_isMin}
  \leanok
  \uses{lem:stationary}
  Let $u$ have the least energy in $u+U_0$, let $W\subseteq u+U_0$ be any set
  and $\tilde u\in W$.  Then $\tilde u$ has the least energy in $W$ if and only
  if $\norm{u-\tilde u}_m\le\norm{u-w}_m$ for every $w\in W$.
\end{lemma}

\begin{proof}
  \leanok
  \uses{lem:stationary}
  On $W$, $\mathcal E(w)=\mathcal E(u)+\tfrac12\norm{w-u}_m^2$ by
  Lemma~\ref{lem:stationary}, and $t\mapsto\sqrt t$ is increasing.
\end{proof}

\begin{lemma}[The Galerkin system]\label{lem:galerkin}
  \lean{Atlas.QuadEnergy.galerkin_iff_isMin}
  \leanok
  \uses{lem:stationary}
  For $w_0\in X$, finitely many trial directions $t_a\in X$ and coordinates
  $c$, write $w(c)=w_0+\sum_ac_at_a$.  Then $w(c)$ has the least energy among
  all $w(c')$ if and only if $m\bigl(w(c),t_a\bigr)=\ell(t_a)$ for every $a$.
\end{lemma}

\begin{proof}
  \leanok
  \uses{lem:stationary}
  Lemma~\ref{lem:stationary} with $U_0$ the span of the $t_a$ and $u=w(c)$.
\end{proof}

\begin{lemma}[The energy size obeys the triangle inequality]\label{lem:energy-triangle}
  \lean{Atlas.QuadEnergy.norm_add_le}
  \leanok
  \uses{def:quad}
  $\norm{x+y}_m\le\norm{x}_m+\norm{y}_m$.
\end{lemma}

\begin{proof}
  \leanok
  The Cauchy--Schwarz inequality $m(x,y)^2\le m(x,x)\,m(y,y)$ for a positive
  semidefinite symmetric form (Mathlib:
  \texttt{LinearMap.BilinForm.apply\_sq\_le\_of\_symm}).
\end{proof}

\begin{definition}[The tier-0 solve]\label{def:superelement}
  \lean{Atlas.Superelement, Atlas.Superelement.Admissible, Atlas.Superelement.U0,
    Atlas.Superelement.quad, Atlas.Superelement.IsSolution,
    Atlas.Superelement.coord, Atlas.Superelement.field,
    Atlas.Superelement.portMatrix, Atlas.Superelement.portLoad,
    Atlas.Superelement.hostMatrix, Atlas.Superelement.hostRhs,
    Atlas.Superelement.fieldError}
  \leanok
  \uses{def:gram, def:assemble, def:quad}
  Pieces $i=1,\dots,P$, each with an energy $M_i$ on $U\times F$
  (Definition~\ref{def:piece}), a source $f_i$, port patterns $q_{i,k}\in F$, a
  matrix $R_i$ reading its port coordinates off the free global coordinates
  $c$, and fixed coordinates $r_i$.  A \emph{global field} $w$ gives every
  piece a pair $w_i\in U\times F$.  Let $\Gamma_0$ be a subspace of families of
  face values (those on which neighbours agree and which vanish on the
  prescribed boundary) and $g$ a family carrying the prescribed boundary
  values.  The field $w$ is \emph{admissible} if its face values minus $g$ lie
  in $\Gamma_0$; $U_0$ is the space of fields with face values in $\Gamma_0$.
  The global energy is
  \[
    \mathcal E(w)=\sum_{i=1}^P\Bigl(\tfrac12M_i(w_i,w_i)-f_i(w_{I,i})\Bigr),
    \qquad\norm{w}_M^2=\sum_{i=1}^PM_i(w_i,w_i),
  \]
  and the \emph{undivided solution} $u$ is an admissible field of least energy
  among admissible fields.  From fields $h_{i,k}$ and $\hat u_{p,i}$ the host
  forms
  \begin{gather*}
    \tilde\Lambda_i=\Lambda_i[h_i],\qquad
    \tilde b_{i,k}=f_i(h_{i,k})-M_i\bigl((h_{i,k},q_{i,k}),(\hat u_{p,i},0)\bigr),\\
    \tilde{\mathsf S}=\sum_iR_i^\top\tilde\Lambda_iR_i,\qquad
    \tilde\chi=\sum_iR_i^\top\bigl(\tilde b_i-\tilde\Lambda_ir_i\bigr),
  \end{gather*}
  and for coordinates $c$ rebuilds, with $x_i=r_i+R_ic$,
  \[
    w(c)_i=\Bigl(\hat u_{p,i}+\sum_kx_{i,k}h_{i,k},\ \ \sum_kx_{i,k}q_{i,k}\Bigr).
  \]
\end{definition}

\begin{definition}[The three assumptions]\label{def:assumptions}
  \lean{Atlas.Superelement.Conforming, Atlas.Superelement.BoundaryDataInPortSpan,
    Atlas.Superelement.EnergyNormIsNorm}
  \leanok
  \uses{def:superelement}
  \begin{enumerate}
    \item \emph{Conforming shared sides:} for every $c$, the family
      $\bigl(\sum_k(R_ic)_kq_{i,k}\bigr)_i$ lies in $\Gamma_0$.
    \item \emph{Boundary data in the port span:} the family
      $\bigl(\sum_kr_{i,k}q_{i,k}\bigr)_i-g$ lies in $\Gamma_0$.
    \item \emph{The energy size is a norm on $U_0$:} $w\in U_0$ and
      $\norm{w}_M=0$ imply $w=0$.
  \end{enumerate}
\end{definition}

\begin{theorem}[The host's system is the Galerkin system]\label{thm:T25-galerkin}
  \lean{Atlas.Superelement.hostSystem_iff_isMin}
  \leanok
  \uses{def:superelement, lem:galerkin}
  For any fields $h,\hat u_p$ and any $c$:
  $\tilde{\mathsf S}c=\tilde\chi$ if and only if
  $\mathcal E(w(c))\le\mathcal E(w(c'))$ for every $c'$.
\end{theorem}

\begin{proof}
  \leanok
  \uses{lem:galerkin}
  The trial directions are
  $t_a=\bigl(\sum_k(R_i)_{ka}(h_{i,k},q_{i,k})\bigr)_i$.  Expanding
  $m(w(c),t_a)=\ell(t_a)$ by bilinearity gives row $a$ of
  $\tilde{\mathsf S}c=\tilde\chi$; then Lemma~\ref{lem:galerkin}.
\end{proof}

\begin{theorem}[T25, energy optimality]\label{thm:T25-optimal}
  \lean{Atlas.Superelement.field_admissible, Atlas.Superelement.energy_optimal}
  \leanok
  \uses{thm:T25-galerkin, lem:cea, def:assumptions}
  Assume 1 and 2 of Definition~\ref{def:assumptions}.  Then every rebuilt
  field $w(c)$ is admissible.  If $u$ is the undivided solution and
  $\tilde{\mathsf S}c=\tilde\chi$, then
  \[
    \norm{u-w(c)}_M\le\norm{u-w(c')}_M\qquad\text{for every }c' .
  \]
\end{theorem}

\begin{proof}
  \leanok
  \uses{thm:T25-galerkin, lem:cea}
  The face values of $w(c)$ minus $g$ are the sum of the two families of
  assumptions 1 and 2.  By Theorem~\ref{thm:T25-galerkin}, $w(c)$ has the least
  energy among the rebuilt fields; Lemma~\ref{lem:cea} with $W$ the set of
  rebuilt fields.
\end{proof}

\begin{theorem}[T25, the error bound]\label{thm:T25-bound}
  \lean{Atlas.Superelement.energy_error_bound}
  \leanok
  \uses{thm:T25-optimal, lem:energy-triangle}
  Under the hypotheses of Theorem~\ref{thm:T25-optimal}, for any reference
  fields $H_{i,k}$, $u_{p,i}$ and any coordinates $c^\star$, with
  $x_i^\star=r_i+R_ic^\star$ and $w^{\mathrm{ref}}$ the field rebuilt from the
  reference fields at $c^\star$,
  \[
    \norm{u-w(c)}_M\ \le\ \norm{u-w^{\mathrm{ref}}}_M
      +\Bigl(\sum_{i=1}^Pa_i(d_i,d_i)\Bigr)^{1/2},
    \qquad d_i=(\hat u_{p,i}-u_{p,i})+\sum_kx^\star_{i,k}\,(h_{i,k}-H_{i,k}).
  \]
  With the exact modes and the classical superelement's coordinates,
  $w^{\mathrm{ref}}=u_m$ and this is the document's bound: truncation error
  plus field errors.
\end{theorem}

\begin{proof}
  \leanok
  \uses{thm:T25-optimal, lem:energy-triangle}
  Let $w^\star$ be the field rebuilt from the network's fields at $c^\star$.
  By Theorem~\ref{thm:T25-optimal} and the triangle inequality,
  $\norm{u-w(c)}_M\le\norm{u-w^\star}_M
  \le\norm{u-w^{\mathrm{ref}}}_M+\norm{w^{\mathrm{ref}}-w^\star}_M$.  The field
  $w^\star-w^{\mathrm{ref}}$ has cell values $d_i$ and zero face values on piece
  $i$, and $M_i\bigl((d,0),(d,0)\bigr)=a_i(d,d)$.
\end{proof}

\begin{theorem}[Solvable for every output; one solution]\label{thm:T25-solvable}
  \lean{Atlas.Superelement.solution_unique, Atlas.Superelement.hostMatrix_posDef,
    Atlas.Superelement.hostSystem_existsUnique}
  \leanok
  \uses{def:assumptions, lem:stationary}
  Assume 3 of Definition~\ref{def:assumptions}.  Then the undivided solution
  is unique.  Assume also 1, that each piece's patterns $q_{i,k}$ are linearly
  independent, and that $R_ic=0$ for every $i$ only for $c=0$.  Then
  $\tilde{\mathsf S}$ is positive definite for every choice of the fields $h$,
  and $\tilde{\mathsf S}c=\tilde\chi$ has exactly one solution.
\end{theorem}

\begin{proof}
  \leanok
  \uses{lem:stationary}
  Two solutions $u,u'$ have equal energy, so $\norm{u'-u}_M=0$ by
  Lemma~\ref{lem:stationary}, and $u'-u\in U_0$.  For the matrix,
  $c^\top\tilde{\mathsf S}c=\norm{t(c)}_M^2$ with
  $t(c)_i=\sum_k(R_ic)_k(h_{i,k},q_{i,k})$, which lies in $U_0$ by assumption
  1.  If it vanishes then $t(c)=0$; its face values $\sum_k(R_ic)_kq_{i,k}$
  vanish, so $R_ic=0$ for every $i$, so $c=0$.
\end{proof}

\section{T24: perturbation of the interface solve}

\paragraph{In plain words.}
The exact interface system is $\mathsf Sc=\chi$ and the host solves a learned
one, $\tilde{\mathsf S}\tilde c=\tilde\chi$.  The two answers differ by at most
the error of the matrix times the size of the exact answer, plus the error of
the load, all divided by the smallest stretch $\beta$ of the learned matrix.
For the Gram construction that constant need not be measured on the learned
matrix: by T25 (iii) the exact problem's own constant serves.

\begin{proposition}[T24]\label{prop:T24}
  \lean{Atlas.interface_perturbation}
  \leanok
  Let $\mathsf S,\tilde{\mathsf S}$ be bounded linear maps between real normed
  spaces with $\beta\norm{x}\le\norm{\tilde{\mathsf S}x}$ for all $x$ and some
  $\beta>0$.  If $\mathsf Sc=\chi$ and $\tilde{\mathsf S}\tilde c=\tilde\chi$,
  \[
    \norm{\tilde c-c}\le
      \frac{\norm{\mathsf S-\tilde{\mathsf S}}\,\norm{c}+\norm{\chi-\tilde\chi}}{\beta}.
  \]
\end{proposition}

\begin{proof}
  \leanok
  $\tilde{\mathsf S}(\tilde c-c)=(\tilde\chi-\chi)+(\mathsf S-\tilde{\mathsf S})c$;
  take norms and use the lower bound.
\end{proof}

\begin{corollary}[T24 with the exact problem's constant]\label{cor:T24-dominates}
  \lean{Atlas.interface_perturbation_of_dominates}
  \leanok
  \uses{prop:T24}
  On a real inner-product space, if
  $\beta\norm{x}^2\le\langle\mathsf Sx,x\rangle\le\langle\tilde{\mathsf S}x,x\rangle$
  for all $x$ and some $\beta>0$, the bound of Proposition~\ref{prop:T24} holds
  with this $\beta$.
\end{corollary}

\begin{proof}
  \leanok
  \uses{prop:T24}
  $\beta\norm{x}^2\le\langle\tilde{\mathsf S}x,x\rangle
  \le\norm{\tilde{\mathsf S}x}\,\norm{x}$ by Cauchy--Schwarz.
\end{proof}

\section{T26: the label-free objective}

\paragraph{In plain words.}
The supervised loss compares the network's fields with exact ones, which a
classical solver must first compute.  The label-free objective needs only the
network's own output and the piece's energy.  The two differ by a number fixed
by the piece and its source, so between any two outputs of the network they
change by the same amount: the same gradients and the same minimisers, with no
solved examples.

\begin{definition}[The two objectives]\label{def:objectives}
  \lean{Atlas.Piece.labelFree, Atlas.Piece.supervised}
  \leanok
  \uses{def:gram}
  For returned fields $h_k$ and $\hat u_p$, a source $f$, exact modes $H_k$ and
  the exact particular field $u_p$ (with $a(v,u_p)=f(v)$ for all $v$):
  \[
    \mathcal L_{\mathrm{en}}=\operatorname{tr}\Lambda[h]+a(\hat u_p,\hat u_p)-2f(\hat u_p),
    \qquad
    \mathcal L_{\mathrm{sup}}=\sum_ka(h_k-H_k,h_k-H_k)+a(\hat u_p-u_p,\hat u_p-u_p).
  \]
\end{definition}

\begin{theorem}[T26]\label{thm:T26}
  \lean{Atlas.Piece.labelFree_eq_supervised_add_const,
    Atlas.Piece.labelFree_sub_eq_supervised_sub}
  \leanok
  \uses{def:objectives, cor:T25-ii}
  $\mathcal L_{\mathrm{en}}=\mathcal L_{\mathrm{sup}}
  +\bigl(\operatorname{tr}\Lambda[H]-a(u_p,u_p)\bigr)$.  The bracket does not
  involve the network's output, so for any two outputs
  $\mathcal L_{\mathrm{en}}-\mathcal L_{\mathrm{en}}'
  =\mathcal L_{\mathrm{sup}}-\mathcal L_{\mathrm{sup}}'$.
\end{theorem}

\begin{proof}
  \leanok
  \uses{cor:T25-ii}
  The trace identity of Corollary~\ref{cor:T25-ii}, and
  $a(\hat u_p-u_p,\hat u_p-u_p)=a(\hat u_p,\hat u_p)-2f(\hat u_p)+a(u_p,u_p)$,
  using $a(\hat u_p,u_p)=f(\hat u_p)$.
\end{proof}
